\label{der:1-3}
记$V_n=\det(x_j^{i-1})$。当$x_j=x_i$时两列相同，$V_n=0$，由多项式因式定理有因子$x_j-x_i$。这些不同的一次因子两两不相伴，在多项式环中为互异素因子，因而其乘积整除$V_n$。排列定义表明每项总次数均为$0+1+\cdots+(n-1)=\frac{n(n-1)}{2}$，与全部差因子的总次数相同，故
\[V_n=C\prod_{i<j}(x_j-x_i).\]
左边单项式$x_1^0x_2^1\cdots x_n^{n-1}$只来自恒等排列，系数为1；右边要取得这一单项式，必须从每个含$x_1$的差中取另一项，再对$x_2$依次处理，最终每个$(x_j-x_i)$均取$x_j$，系数也是1，所以$C=1$。

对秩一更新，将第$j$列写为$e_j+v_j\vu$。展开后，若选两个更新列，它们均平行于$\vu$，该项为零。不选更新列得1；只更新第$j$列时，沿该列展开或用标准基位置可得行列式$u_j$，系数为$v_j$。因此
\[\det(I+\vu\vv\T)=1+\sum_jv_ju_j=1+\vv\T\vu.\]
对$D+\vu\vv\T$，不选更新列给$\prod_i d_i$；只更新第$j$列时，其他列固定为$d_ie_i$，提出$\prod_{i\ne j}d_i$后只剩$u_jv_j$。所以无条件成立
\[\det(D+\vu\vv\T)=\prod_i d_i+\sum_j u_jv_j\prod_{i\ne j}d_i.\]
当全部$d_i\ne0$，提出$\prod_i d_i$即得含分式版本。这里无除法推导已经包含$d_i=0$的情况。
